\begin{examplebox}
\textbf{例 1：焦点和准线已知。} $F(2,0)$，$\ell:x=-2$。
顶点为 $(0,0)$，焦准距 $p=4$，直接得到 $\boxed{y^2=8x}$。
也可从 $\sqrt{(x-2)^2+y^2}=|x+2|$ 一次平方得到。
回验：点 $(2,4)$ 两个距离均为 $4$。

\medskip
\textbf{例 2：焦点换到竖直方向。} $F(0,1)$，$\ell:y=-1$。
$p=2$，开口向上，故 $\boxed{x^2=4y}$。
回验：点 $(2,1)$ 到焦点距离为 $2$，到准线距离也为 $2$。
\end{examplebox}
\clearpage
\begin{examplebox}
\textbf{例 3：最小值，先把焦半径换成垂距。}
$P$ 在 $y^2=4x$ 上，$A(3,2)$，求 $PA+PF$ 的最小值。
\begin{center}\begin{tikzpicture}[font=\small,>=stealth]
\draw[black!30,rounded corners=2pt] (0,0) rectangle (14.4,8.2);
\node[anchor=west,font=\small\bfseries] at (.5,7.72) {例 3：用 $PF=PH$ 把两段距离拉成一条水平线};
\begin{scope}[shift={(4.2,2.35)},scale=.9]
\draw[violet,thick] (-1,-.5)--(-1,4.6) node[above] {$\ell:x=-1$};
\draw[blue,very thick,domain=-.5:4.5,samples=121,variable=\t] plot ({\t*\t/4},\t);
\draw[black!45] (-1.5,0)--(5.5,0);
\draw[orange!85!black,very thick] (-1,2)--(3,2);
\draw[violet,dashed,thick] (3,2)--(4,4)--(-1,4);
\draw[black!45,dashed] (1,0)--(1,2);
\fill (1,0) circle (2pt) node[below] {$F$};
\fill (1,2) circle (2pt) node[above left=6pt] {$P(1,2)$};
\fill (3,2) circle (2pt) node[right] {$A(3,2)$};
\fill (-1,2) circle (2pt) node[left] {$H$};
\fill (4,4) circle (2pt) node[below right=8pt] {$Q$};
\fill (-1,4) circle (2pt) node[left] {$H_Q$};
\end{scope}

\node[anchor=west] at (.7,.55) {$AP+PF=AP+PH\ge d(A,\ell)=4$；$P=(1,2)$ 时等号成立};
\end{tikzpicture}\end{center}

设 $H$ 为 $P$ 的准线垂足。由定义，$PA+PF=PA+PH$。
折线 $A\to P\to H$ 不短于 $A$ 到准线的垂距，故总长至少为 $4$。
取 $P=(1,2)$，它既在抛物线上，又在线段 $AH$ 内，
且 $PA=2,PF=2$，所以最小值为 $\boxed{4}$，等号可实现。

\textbf{代数角度回验。} $x\ge0$，$PF=x+1$，而 $PA\ge|3-x|$。
因此 $PA+PF\ge |3-x|+x+1\ge4$；$P=(1,2)$ 达到等号。
推广到其他定点时，必须另查垂足连线上的交点是否落在相应线段内，不能直接照搬取等条件。
\end{examplebox}
\clearpage

% ===== C038 补充：原三个例题与原有 TikZ 完整保留 =====
\begin{examplebox}
\textbf{例 4：考场反向识别焦点与准线（不重新推导）。}
给出 $y^2=4x$，求焦点、准线及最小焦半径。

\textbf{第一眼：}$2p=4$，所以 $p=2$。\textbf{第二眼：}开口向右，
焦点 $F=(1,0)$，准线 $\ell:x=-1$。
\textbf{第三眼：}顶点处的焦半径最短，为 $p/2=1$。

核对 $P=(1,2)$：
\[
PF=\sqrt{(1-1)^2+2^2}=2,\qquad d(P,\ell)=|1+1|=2.
\]
\textbf{只要 2 步：看系数求 $p$，再把 $p$ 平分到顶点两侧。}
\end{examplebox}
\clearpage
\begin{examplebox}
\textbf{例 5：顶点不在原点时，也可从“中点”立即写出方程。}
已知焦点 $F(1,2)$，准线 $\ell:x=-3$，求轨迹方程。

准线上的垂足为 $D(-3,2)$，焦准距 $FD=4$，中点（顶点）
$O=(-1,2)$。焦点在准线右边，所以开口向右。
把 $y^2=2px$ 中的 $x$ 换为 $x+1$，$y$ 换为 $y-2$：
\[
\boxed{(y-2)^2=8(x+1)}.
\]
\textbf{数值检验：}取 $P=(1,6)$，左边为 $16$，右边也为 $16$；
其到焦点的距离为 $4$，到准线的距离也为 $4$。

\textbf{考场口令：}先找顶点（焦点和垂足的中点），再看开口，最后平移标准式。
\end{examplebox}
\clearpage
\begin{examplebox}
\textbf{例 6：负号出现时，先判断开口方向。}
\begin{enumerate}[leftmargin=*]
\item $y^2=-12x$。它向左开口，$2p=12$，所以 $p=6$。
焦点 $F=(-3,0)$，准线 $\ell:x=3$。
\item $x^2=6y$。它向上开口，$2p=6$，所以 $p=3$。
焦点 $F=(0,3/2)$，准线 $\ell:y=-3/2$。
\end{enumerate}
\textbf{口令：}平方的是哪个坐标，轴就沿\textbf{另一个坐标轴}；
一次项的正负号，决定沿该轴的开口方向。

\medskip
\textbf{例 7：一个点是否在抛物线上？}
仍取 $y^2=4x$。对 $Q=(4,4)$，$y^2=16=4x$，故在抛物线上；
再用定义回验，$QF=5$，$d(Q,\ell)=5$。
但 $R=(4,3)$ 满足 $9\ne16$，所以不在抛物线上。
\textbf{注意：}用一个数值点只能检验这个点，不能替代整条轨迹的证明。
\end{examplebox}
\clearpage
\begin{examplebox}
\textbf{例 8：最值技巧什么时候不能直接取等？}
% Equal units; the intersection of y=2 and y^2=4x is always P=(1,2).
\begin{center}\begin{tikzpicture}[font=\small,>=stealth]
\draw[black!30,rounded corners=2pt] (0,0) rectangle (14.4,8.2);
\node[anchor=west,font=\small\bfseries] at (.5,7.72) {同一个交点 $P$：在线段内可以取等，在线段外不能取等};
\node at (3.6,6.65) {$A=(3,2)$：$P\in\overline{AH}$};
\node at (10.8,6.65) {$A=(0,2)$：$P\notin\overline{AH}$};
\foreach \cx/\ax in {3.6/3,10.8/0} {
  \begin{scope}[shift={(\cx,3.7)},scale=.7]
    \draw[black!40] (-1.5,0)--(4.3,0);
    \draw[violet,thick] (-1,-.5)--(-1,3.4);
    \node[below left=3pt] at (-1,-.5) {$\ell$};
    \draw[blue,very thick,domain=-.5:3.4,samples=101,variable=\t]
      plot ({\t*\t/4},\t);
    \draw[black!45,dashed] (-1.4,2)--(4,2);
    \draw[violet,very thick] (-1,2)--(\ax,2);
    \draw[orange!85!black,thick] (1,0)--(1,2);
    \draw[black!60] (-1,1.75)--(-.75,1.75)--(-.75,2);
    \fill (1,0) circle (2pt) node[below=5pt] {$F$};
    \fill (-1,2) circle (2pt) node[left=5pt] {$H$};
    \fill (1,2) circle (2pt) node[below right=7pt] {$P$};
  \end{scope}
}
\fill (5.7,5.1) circle (1.4pt) node[right=6pt] {$A$};
\fill (10.8,5.1) circle (1.4pt) node[above=7pt] {$A$};
\node at (3.6,2.25) {$AP+PF=2+2=4$};
\node at (10.8,2.25) {$P=(1,2)$ 不在所需线段上};
\node at (3.6,1.4) {下界 $d(A,\ell)=4$ 可取};
\node at (10.8,1.4) {下界 $d(A,\ell)=1$ 不可取};
\node[anchor=west] at (.7,.55) {两图同尺度；紫色为线段 $AH$，取等要求 $P$ 在线段 $AH$ 上，且 $AH\perp\ell$};
\end{tikzpicture}\end{center}

在 $y^2=4x$ 上取 $P$，另给 $A=(0,2)$。
定义仍可得到
\[
PA+PF=PA+PH\ge d(A,\ell)=1.
\]
但不能就说最小值为 $1$。若要取等，$A,P,H$ 必须共线，
且 $P$ 在 $A$ 与准线垂足 $H=(-1,2)$ 之间。
这些点的横坐标满足 $-1\le x\le0$。
可是抛物线与 $y=2$ 的交点只有 $P=(1,2)$，\textbf{不在该线段上}。

因此这个下界虽成立，\textbf{在本题却取不到等号}，实际最小值严格大于 $1$。
图中 $P=(1,2)$ 只用来检查取等位置，并非断言它给出右图题设的最小值。
这说明：\textbf{距离转化先给下界，再检验取等点}；两步缺一不可。
\end{examplebox}
